\documentclass[12pt,a4paper,onecolumn]{article}

% Encoding and language
\usepackage[utf8]{inputenc}
\usepackage[T1]{fontenc}
\usepackage[english]{babel}

% Font
\usepackage{times}

% Set margins
\usepackage[a4paper, top=2cm, bottom=2cm, left=2.5cm, right=2.5cm]{geometry}

% Tabulation
\usepackage{tabto}

% Line spacing
\usepackage{setspace}
\onehalfspacing

% Page numbers
\usepackage{fancyhdr}
\pagestyle{fancy}
\fancyhf{} % Clear default header/footer
\rfoot{\thepage} % Right-aligned page numbers in footer

% Remove page number from cover/abstract/etc
\usepackage{titlesec}
\newcommand{\noPageNumber}{\thispagestyle{empty}}

% Prevent indentation
\setlength{\parindent}{0pt}
\setlength{\parskip}{0.5em}

% Figure management
\usepackage{graphicx}
\usepackage{caption}
\usepackage{subcaption}

% Table
\usepackage{booktabs, tabularx}
\renewcommand\tabularxcolumn[1]{m{#1}}

\usepackage{pgfplots}
\pgfplotsset{compat=1.16}
\def\mathdefault#1{#1}
\everymath=\expandafter{\the\everymath\displaystyle}

% Subfigure management
\usepackage{layouts}
%\usepackage{floatrow}
%\usepackage[label font=bf,labelformat=simple]{subfig}% <-- changed
%\floatsetup[figure]{style=plain,subcapbesideposition=top}


% Optional: headings formatting
\titleformat{\section}{\large\bfseries}{\thesection}{1em}{}
\titleformat{\subsection}{\normalsize\bfseries}{\thesubsection}{1em}{}

% Optional: remove widow/orphan lines
\widowpenalty=10000
\clubpenalty=10000

% Mathematics
\usepackage{amsmath,mathtools,amssymb} 

\begin{document}
 %% textwidth in cm: \printinunitsof{in}\prntlen{\textwidth}

\begin{table}[htb]
\centering
\footnotesize
\begin{tabularx}{\linewidth}{c >{\raggedright\arraybackslash}X c c c}
\toprule
Parameter & Definition & Unit & Par A & Par B \\

\midrule
\textit{n} & Number of host LCT subcompartments. Equals to $1/CV_H^2$ with $CV_H$ the
coefficient of variation of host maturation time & unit-less & \textit{varies} & \textit{varies} \\

\textit{m} & Number of parasitoid LCT subcompartments. Equals to $1/CV_P^2$ with $CV_P$ the
coefficient of variation of parasitoid maturation time & unit-less & \textit{varies} & \textit{varies} \\

$m_{H1}$ & Maturation rate of H1. Average maturation time  $T_{H1}$, can be defined as $1/m_{H1}$ & $t^{-1}$ & 1  & 1 \\

$\delta_{H1}$ & Background mortality of H1 & $t^{-1}$ & 0 & 0.00425 \\
$\delta_{H2}$ & Background mortality of H2. Host adult longevity, $T_{H2}$ can be defined as $1/\delta_{H2}$ & $t^{-1}$ &  0.2 & 0.2  \\
$^\rho$ & Average number of offspring of one adult during its lifetime. Equals to $r/\delta_{H2}$ with r the adult fecundity rate. & unit-less & 33 & 105 \\

$m_{P1}$ & Maturation rate of P1. Average maturation time  $T_{P1}$, can be defined as $1/m_{P1}$ & $t^{-1}$ & 1 & 0.05 \\
$\delta_{P1}$ & Background mortality of P1 & $t^{-1}$ & 1 & 0.1  \\
$\delta_{P2}$ & Background mortality of P2 & $t^{-1}$ & 8 & 0.5  \\
$a$ & Rate of host discovery by a single parasitoid & $t^{-1}$ & 1 & 0.03  \\

$\theta_H$ & Rate parameter of Erlang maturation time distribution. Set equal to $n \cdot m_{H1}$ (Appendix A.1.3) & $t^{-1}$ & \textit{varies}  & \textit{varies}   \\




\bottomrule
\end{tabularx}
\caption{Parameters and values of Model I (Murdoch 1987). Par A : parameter set inspired from those in Murdoch (1987). Par B : parameter set inspired from measurements in Sait et al. (1994) and Harvey et al. (1994) on the \textit{Plodia-Venturia} host-parasitoid system.}
\end{table}





















\begin{table}[htb]
\centering
\footnotesize
\begin{tabularx}{\linewidth}{c >{\raggedright\arraybackslash}X c c c}
\toprule
Parameter & Definition & Unit & Exploratory value \\

\midrule
$c_S$ & Colonization-Establishment-Growth-through-adulthood rate of shrubs (herbivore-indep). If CC=1, shrubs colonize up to $c_S$ times their initial (low) cover in 1 time step (year). & per shrub occupation, per open space, $t^{-1}$ & 1 \\
$c_T$ & Colonization-Establishment of tree seedling-sapling (herbivore-indep)  & per adult tree occupation, per open space, $t^{-1}$ & 0.2 \\
$c_{S,T}$ & Colonization-Establishment-Growth-through-adulthood rate of shrubs in a sapling-occupied locality (herbivore-indep) & per adult tree occupation, per open space, $t^{-1}$ & $c_S$ or 0 (no T1-N link) \\
$c_{T,S}$ & Colonization-Establishment of tree seedling-sapling in a shrub-occupied locality (herbivore-indep) & per adult tree occupation, per open space, $t^{-1}$ & $c_T$ \\

$g_1$ & Proportion of T1 patch that become invulnerable to herbivory per time step & $t^{-1}$ per T1 occupation & 0.2  \\
$g_{1,S}$ & Proportion of (sapling co-occuring with shrubs) patch that become invulnerable to herbivory per time step  & $t^{-1}$ per N=(T1,S) occupation & $g_1$ (no competition) \\
$m_{T2}$ & Extinction (proportion of patch) rate of T2 & $t^{-1}$ per T1 occupation & 0.02  \\

$a_S$ & Proportion of shrub patches that are cleared by a herbivory pressure of 1 each time step & per shrub occupation, per herbivore density, $t^{-1}$ & 0.1 \\
$a_T$ & Proportion of seedling-sapling patches that are cleared by a herbivory pressure of 1 each time step  & per sapling occupation, per herbivore density, $t^{-1}$ & 1 \\
$a_{S,T}$ & Per capita attack (= perturbation that causes extinction) rate of shrubs co-occurying with tree saplings & per sapling occupation, per herbivore density, $t^{-1}$ & $a_S$ or 0 (no T1-N link) \\
$a_{T,S}$ & Per capita attack (= perturbation that causes extinction) rate of saplings co-occurying with shrubs & per sapling occupation, per herbivore density, $t^{-1}$ & $a_T$ (no AR) or $a_S$ (AR) \\

$CC$ & Total occupation of all vegetation patches caused by finite space & patch proportion & 1 \\

\bottomrule
\end{tabularx}
\caption{Parameters, meaning and values of first Vera model}
\end{table}



\end{document}